\chapter*{Podsumowanie}
\addcontentsline{toc}{chapter}{Podsumowanie}

Wynikiem pracy jest aplikacja, której funkcją jest implementacja graficznego interfejsu do pobierania danych tekstowych z serwisu YouTube, przedstawienie narzędzi do ustawienia parametrów i wyboru metod do analizy danych tekstowych oraz końcowy wykres odzwierciedlający dane tekstowe.\\

Aplikacja pozwala na efektywne pobieranie danych tekstowych z komentarzy w serwisie YouTube, co umożliwia badanie opinii i nastrojów użytkowników. Dzięki zastosowaniu interfejsu graficznego użytkownicy mogą łatwo ustawiać parametry analizy oraz wybierać różne metody przetwarzania danych tekstowych, takie jak analiza sentymentu, identyfikacja kluczowych słów czy grupowanie danych.\\

Podsumowując powyższe wyniki, można podkreślić kilka kluczowych wniosków. Aplikacja jest skutecznym narzędziem do analizy danych tekstowych pochodzących z serwisu YouTube. Pozwala na szybkie pobieranie i przetwarzanie dużej ilości danych, co może być szczególnie użyteczne w~~badaniach opinii publicznej. Aplikacja oferuje różne metody przetwarzania danych, co pozwala na wszechstronną i obiektywną ocenę danych tekstowych. Wykorzystano metody klasteryzacji, takie jak K-means i DBSCAN oraz metody wektoryzacji, takie jak TF-IDF i Word2Vec. Dodatkowo, zastosowano techniki redukcji wymiarowości, takie jak t-SNE. Użytkownicy mogą wybierać spośród wymienionych algorytmów i technik, co zwiększa elastyczność i dokładność analizy. Graficzny interfejs użytkownika jest intuicyjny i łatwy w obsłudze. Dane na wykresie są przedstawione za~~pomocą klastrów o różnych kolorach, co ułatwia szybkie zorientowanie się wśród różnych klastrów na wykresie. Kolorowe klastry pozwalają na łatwe rozróżnienie i interpretację wyników analizy.\\

Pracę można rozwinąć poprzez integrację z innymi platformami społecznościowymi, takimi jak Twitter czy Facebook oraz wdrożenie bardziej zaawansowanych algorytmów analizy tekstu, takich jak modele głębokiego uczenia. Możliwe jest także dodanie funkcji automatycznego doboru optymalnych parametrów analizy danych.\\